\chapter{Conclusiones}
\label{cap:concl}

\section{¿Qué metaheurística funcionó mejor?}
La que mejor se adaptó a los grafos difíciles del FHCP fue la \textbf{colonia artificial de abejas
(ABC)}. Tuvo la tasa de éxito más alta en el experimento A (empatada con \alg{ACO + Pósa}),
en el benchmarking resolvió más del doble de grafos que \alg{ACO + Pósa}
(\resueltasBABC{} de \grafosB{}, contra \resueltasBACOposa) y fue la que quedó más cerca de la solución en los grafos grandes,
con menos del 1\,\% de aristas faltantes incluso con 9\,528 vértices.

La \textbf{colonia de hormigas con rotaciones de Pósa} fue igual de buena en los grafos pequeños,
y a veces más rápida, pero no escaló: a partir de unos 200 vértices dejó de resolver grafos.

\section{Lo que aprendimos}
\begin{enumerate}
  \item \textbf{La representación y los movimientos importan más que el algoritmo.} Los tres
        algoritmos con éxito usan información del grafo para moverse: ABC repara conexiones con
        aristas reales, ACO camina solo por aristas que existen y usa rotaciones, y SA (mixto)
        invierte tramos. Los que se mueven a ciegas (GA, SA con intercambios) o sobre números
        reales (DE y PSO) no encontraron ningún ciclo.
  \item \textbf{Un cambio pequeño puede hacer una gran diferencia.} SA solo con intercambios no
        resolvió nada, pero con la mitad de los movimientos cambiados a inversiones logró
        \exitosASA{} éxitos con el mismo presupuesto. Lo mismo pasó con ACO: sin rotaciones, 0
        éxitos; con rotaciones, \exitosAACOposa.
  \item \textbf{Las random keys no funcionan bien para este problema.} Garantizan rutas válidas,
        pero un cambio pequeño en las claves puede reordenar media ruta, así que DE y PSO casi no
        pueden mejorar poco a poco. En los grafos grandes quedaron prácticamente igual que una
        ruta al azar.
  \item \textbf{Los grafos cúbicos son el gran cuello de botella.} Casi ninguno se resolvió
        (solo 1 de los 42 grafos cúbicos que hay entre los primeros 100). Con solo tres vecinos por vértice, cualquier error
        deja vértices sin salida.
  \item \textbf{No encontrar el ciclo no significa que no exista.} Todos los grafos del FHCP
        tienen ciclo hamiltoniano, así que cada corrida sin éxito es un límite de nuestros
        algoritmos, no del grafo.
\end{enumerate}

\section{Limitaciones}
\begin{itemize}
  \item La comparación fue a igual tiempo y no a igual número de evaluaciones
        (tabla~\ref{tab:presupuesto}), así que algunos algoritmos usaron muchas más evaluaciones
        que otros.
  \item Las rotaciones de Pósa no se cuentan como evaluaciones, lo que hace que \alg{ACO + Pósa}
        parezca más barato de lo que es.
  \item Usamos los parámetros iniciales de cada algoritmo, sin una calibración sistemática sobre el
        FHCP. Es posible que GA, DE o PSO mejoren con otros valores.
  \item En el experimento B usamos 2 semillas por grafo y en el C solo una, así que esos resultados
        son menos seguros que los del experimento A.
  \item El código está en Python puro, así que los tiempos no se pueden comparar directamente con
        los de herramientas escritas en C o C++.
\end{itemize}

\section{Trabajo futuro}
\begin{itemize}
  \item Calibrar los parámetros de cada algoritmo con algunos grafos reservados para eso.
  \item Agregar a DE y PSO un paso de búsqueda local con la inversión guiada, para ver si así
        logran competir.
  \item Dar más tiempo a los grafos que no se resolvieron, en especial a los cúbicos.
  \item Combinar lo mejor de ABC y de ACO + Pósa: construir rutas con hormigas y luego mejorarlas
        con abejas.
\end{itemize}
